En el grado noveno, los estudiantes del curso están organizando una venta de postre de fresa en la clase de emprendimiento. Para ello, utilizan vasos que tienen forma prácticamente cilíndrica, con un radio de 3 cm y una altura de 8 cm.

Saben que una caja de postre de fresa contiene 2 litros (2000 cm$^3$) de producto y desean estimar cuántos helados pueden servir para planear la venta.

\begin{enumerate}
\item Realice un dibujo que represente el vaso con sus dimensiones.
\item Calcule el volumen de un vaso de helado.
\item Determine aproximadamente cuántos vasos de helado se pueden obtener con una caja.
\end{enumerate}

Sustente sus procedimientos con argumentos matemáticos claros. No se aceptan respuestas sin justificación.